\section{Graphs, Contact Fields, and Physical Grounding}
\label{sec:structured}

\subsection{A graph is a hypothesis about factorization}

Interaction Networks established a general relational inductive bias: entities
and their pairwise relations can be modelled separately, with messages aggregated
to update entity states~\cite{battaglia2016interaction}. NerveNet applies graph
policies over articulated morphology~\cite{wang2018nervenet}; graph simulators
learn physical dynamics over interacting particles or bodies
~\cite{sanchez2020gns}; Graphormer demonstrates that Transformer attention can be
modified with graph-structural encodings~\cite{ying2021graphormer}; and AnyMorph
studies policies across morphologies~\cite{trabucco2022anymorph}. Together these
works make ``uses a GNN'' an architectural choice, not a sufficient novelty claim.

For tactile bimanual manipulation, a graph encodes a stronger hypothesis: task
dynamics factorize around interfaces and the paths connecting them. This
hypothesis is useful only if (i) node and edge semantics correspond to deployable
or inferable quantities, (ii) the factorization improves sample efficiency,
robustness, interpretability, or intervention response, and (iii) the same gain
cannot be explained by additional state or parameters.

\subsection{Taxel, finger, and morphology graphs}

Graph construction at the tactile-array or finger level is relatively direct.
Distributed tactile sensors have known positions on an articulated hand;
convolution or message passing can share information along this geometry.
Multi-finger in-hand manipulation with distributed tactile GCNs demonstrates this
idea~\cite{funabashi2022distributed}. TacGNN uses hierarchical graph structure for
blind tactile-based in-hand manipulation~\cite{yang2023tacgnn}. TactiGraph treats
event-based tactile measurements asynchronously for contact-angle prediction
~\cite{sajwani2023tactigraph}, and TouchWGNN develops spatio-temporal tactile
graphs for multimodal dexterous manipulation~\cite{ning2026touchwgnn}.

These models address meaningful spatial and temporal organization. However, a
taxel adjacency graph is not the same as a dynamic task contact graph. Taxels can
remain neighbours in the sensor layout whether or not their links contact the
same object; fingers can be kinematically adjacent while transmitting loads
through different entities. Bimanual load-path work therefore needs at least two
levels: relatively stable morphology/sensor structure and dynamic inter-entity
interfaces.

\subsection{Dynamics graphs and contact discontinuities}

RoboPack learns tactile-informed recurrent graph dynamics for dense packing and
uses the model for planning/control~\cite{ai2024robopack}. Its relational
factorization is compelling because tool, objects, and environment interact over
time. ContactNets targets discontinuous contact dynamics with a structured
implicit model~\cite{pfrommer2020contactnets}; graph network simulators have been
shown to model rigid contact transitions~\cite{allen2022contact}. This line of
work suggests two design principles for bimanual representation: predict changes
at interfaces, and avoid treating discontinuity as ordinary smooth motion noise.

MeshPriorDiT adds a useful hierarchical counterexample to an architecture-only
comparison. An action-conditioned mesh GNN supplies topology and grasp-constraint
priors, while a flow-matching diffusion Transformer models residual and global
cloth effects~\cite{wang2026meshpriordit}. Its evidence concerns deformable mesh
dynamics rather than tactile load transfer, but the factorization is directly
relevant: structured local priors and expressive global residuals can be
complementary. A bimanual study should therefore compare graph-only,
global-residual-only, and combined models rather than forcing a false GNN-versus-
Transformer dichotomy.

Yet dynamics prediction can be correct for the wrong reason. Exact object pose,
collision geometry, and contact labels can make a simulator graph accurate while
an observation-only system remains underdetermined. A useful paper must specify
which signals construct the graph at training and at deployment. ``Physics
grounded'' should describe targets, conservation constraints, frames, or
interventions---not merely that the task occurs in a physics engine.

Dex-X makes the distinction between removing privileged state and removing
reference dependence concrete~\cite{chen2026dexx}. Its spatialized tactile-force
student drops current object-pose and geometric privileged blocks but retains
next-step motion references and the demonstration's target object pose. Its
bimanual handover result belongs to simulation; real evaluation includes
single-hand manipulation with strongly mixed object-generalization outcomes.
This is a useful spatialized-force and distillation comparator, but should not be
called reference-free, real bimanual load-path inference, or topology transfer.
CoRL acceptance was author-reported at this audit, not independently confirmed
from an official venue record.

\subsection{Contact fields, modes, and anchors}

Graphs are not the only structured representation. Neural Contact Fields track
extrinsic contact through tactile sensing~\cite{higuera2023ncf}; object-centric
contact-field approaches extend this spatial view to visuo-tactile state
estimation and tool interaction. Semantic Contact Fields predict semantically
meaningful contact/force structure for tool manipulation
~\cite{ma2026semanticcontact}. Contact-mode control represents discrete
constraint regimes for non-prehensile manipulation~\cite{oller2024extrinsic}.
Contact-anchored policies and contact-grounded policy learning similarly use
contact as a privileged or learned organizing variable
~\cite{cui2026contactanchored,xu2026contactgrounded}.

Recent work also clarifies two meanings of ``contact structure.'' One-shot
learning by identifying physical interactions extracts making/breaking contact
events from a demonstration and uses them as transition conditions for hybrid
position--force control~\cite{overbeek2026oneshot}. CoToGrasp conditions grasp
synthesis on a requested contact topology in a canonical workspace
~\cite{merand2026cotograsp}. Both make structure explicit, but the former uses a
demonstrated event state machine and the latter a prescribed grasp-contact layout.
Neither is equivalent to inferring a time-varying cross-hand topology from raw
deployment observations. They are valuable controls precisely because they test
how much structure can be specified rather than learned.

CADeT is positive evidence that useful interaction structure need not be supplied
as simulator truth~\cite{hu2026cadet}. It infers a belief over predefined coupled,
decoupled, and blocked tissue-transmission modes from visually tracked responses
and physical probes. In 150 simulation trials with mode truth withheld, it
identifies 149 modes correctly; matched information-guided probes require fewer
additional actions than random probes. Phantom and ex-vivo experiments extend
the evidence beyond simulation. This is visual, fixed-family mode inference,
with nominal calibration before blockage, rather than tactile cross-hand graph
inference or held-out-participant topology transfer. GP covariance informs mode
likelihoods but is not propagated through the MPC horizon; entropy convergence
is not probability-calibration evidence, and contact force/strain safety is not
measured or guaranteed.

Fields are attractive when contact location varies continuously over object
surfaces; discrete graphs are attractive when entity/interface identity and
message routing matter. Hybrid representations can use a field to infer contact
location and confidence, then instantiate an entity-level relation. The decisive
comparison should be empirical: a field, graph, anchored token set, and
unstructured attention model should receive the same observations and predictive
targets.

\begin{table}[t]
\centering
\caption{Structured-representation choices. The final column states what must be
demonstrated rather than assumed.}
\label{tab:structured_models}
\begin{tabularx}{\linewidth}{P{0.18\linewidth} P{0.21\linewidth} P{0.25\linewidth} Y}
\toprule
\textbf{Structure} & \textbf{Natural unit} & \textbf{Typical source} & \textbf{Validity test} \\
\midrule
Sensor/morphology graph & taxel, link, joint & calibrated layout and kinematics & robustness across layouts, missing sensors, and embodiments \\
Static task graph & hand, object, tool & designer-specified entity relations & advantage over matched attention without dynamic edge changes \\
Dynamic contact graph & active interface & tactile/proprioceptive inference or simulator contact & inferred--oracle gap, transition timing, false-edge interventions \\
Object-centric anchor & local feature in object frame & visual/kinematic pose estimate & degradation under pose error and unseen grasps \\
Contact field & continuous surface point/force & tactile/vision and geometry & contact localization, force consistency, surface/OOD generalization \\
Contact mode & discrete constraint regime & estimator, demonstration, or privileged label & mode-transition accuracy and closed-loop recovery \\
Physics residual/constraint & force balance, momentum, complementarity & measured or simulated quantities & calibration, counterfactual consistency, no label circularity \\
\bottomrule
\end{tabularx}
\end{table}

\subsection{PhysGraph as a boundary case}

PhysGraph is a directly relevant early graph-policy formulation for bimanual
hand--tool--object manipulation~\cite{li2026physgraph}. It organizes hand, tool,
and object information in a graph Transformer and reports strong simulated
control. Its input regime, however, includes simulator state, fingertip resultant
forces, object/tool geometric encodings, and future reference. It should therefore
be read as evidence for relational policy factorization under privileged state,
not as tactile-array representation pretraining or observation-grounded contact
inference.

A local audit of the public main-branch snapshot at the corpus freeze found that
the bias generator exposes contact-mask and geometric-position inputs, while the
default forward call does not pass them; the resulting public path uses the
configured topology rather than the advertised dynamic contact/geometric bias.
This is a reproducibility observation about that code snapshot, not evidence that
the authors' reported experiments used the same path. The correct experimental
response is to reproduce both the stock public path and a clearly documented
corrected-bias path, then avoid attributing their difference to the paper without
author confirmation.

This case illustrates why code-level observability matters. A diagram can show
dynamic edges while the executed model uses static connectivity; a paper can
describe tactile feedback while the input is a low-dimensional simulator force;
and a physically plausible target can be constructed from unavailable future
state. Reproduction should trace every tensor from sensor or simulator source to
the model and record its deployment availability.

\subsection{Physical objectives for cross-interface structure}

Suppose an inferred interface edge $e_{ij,t}$ has probability $p_{ij,t}$ and an
estimated wrench $\hat{\bm w}_{ij,t}$ expressed in a common object frame. Useful
objectives can combine:

\begin{equation}
\mathcal{L} = \lambda_e \mathcal{L}_{\text{edge}}
 + \lambda_w \mathcal{L}_{\text{future-wrench}}
 + \lambda_s \mathcal{L}_{\text{slip}}
 + \lambda_p \mathcal{L}_{\text{phase}}
 + \lambda_c \mathcal{L}_{\text{consistency}}.
\label{eq:physical_loss}
\end{equation}

The consistency term might penalize disagreement between predicted interface
wrenches and measured object acceleration or enforce action--reaction relations
when both sides of an interface are observed. It should be soft and
uncertainty-aware: unmodelled compliance, gravity-compensation error, and
unobserved environmental contact make exact balance inappropriate. The model
should also predict confidence, because a missing edge under sensor occlusion is
not equivalent to a confidently absent contact.

For sustained contact, the state should additionally represent slow relaxation.
Temperature-coupled viscoelastic force estimation shows that identical-looking
instantaneous configurations can correspond to different force histories
~\cite{wang2026relaxation}. Interface objectives should therefore be evaluated at
multiple contact durations and, where relevant, condition on temperature or an
identified relaxation state rather than treating drift as independent noise.

The central empirical question remains comparative. If a temporal Transformer
with the same inputs predicts future wrench and phase equally well, explicit
edges may add interpretability but not predictive value. If structure improves
data efficiency, transition accuracy, or transfer to a new topology---and the
benefit disappears under random or time-shuffled edges---the relational
hypothesis receives stronger support.
